% Image credits for the photographs and AI illustrations of Book 2 (High School), Arabic edition.
% Machine-readable license records live in images/book2/CREDITS.md; keep this
% file in sync with frontmatter/image-credits-book2.tex (one \item per item there).
% As in the Book 1 Arabic credits: personal names are transliterated, with the
% original spelling in parentheses (the credit itself); institutions, datasets,
% works and licences stay verbatim inside \babelsublr; only the gloss is translated.
\chapter*{مصادر الصور}

الرسوم والخرائط والمخططات في هذا الكتاب أصلية؛ وما رُسم منها من بيانات
منشورة يذكر مصدره في التعليق المرافق أو في هذه القائمة. أما الصور
الفوتوغرافية فمستعملة بمقتضى رخصها الحرة، مع الشكر لأصحابها؛ وروابط المصادر
الكاملة وروابط الرخص لكل صورة مدرجة في \texttt{images/book2/CREDITS.md} في
مستودع الكتاب.

\bigskip
وإلى جانب الصور الفوتوغرافية والرسوم الأصلية، بعض الأشكال رسوم توضيحية
مولَّدة بالذكاء الاصطناعي، أُنتجت بأداة توليد الصور من OpenAI انطلاقًا من
أوامر كتبها محررو الكتاب، ورُوجعت دقتها العلمية قبل إدراجها. وهي لا تُظهر
إلا مشاهد عامة، ولا تمثّل أبدًا مكانًا حقيقيًا أو شخصًا حقيقيًا أو عيّنة
حقيقية. والأمر الكامل لكل صورة مولَّدة مسجَّل في
\texttt{images/book2/ai/PROMPTS.md} في المستودع.

\bigskip
\noindent الصور الفوتوغرافية والبيانات، فصلًا فصلًا:
\begin{itemize}\small
  \item الفصل~1 (قياس الأرض): الأرض كما رُئيت من \babelsublr{Apollo~17}؛ صورة من \babelsublr{NASA}، ملك عام (عبر \babelsublr{Wikimedia Commons}).
  \item الفصل~1: دائرة التكرار، لوحة من كتاب ج.-ب.~دولامبر (\babelsublr{J.-B. Delambre})، \babelsublr{\emph{Base du système métrique décimal}}، المجلد~2 (1807)؛ ملك عام (عبر \babelsublr{Wikimedia Commons}، من \babelsublr{BnF Gallica}).
  \item الفصل~1: أطوال الدرجات محسوبة على المجسم الإهليلجي \babelsublr{WGS~84} (مجموعة بيانات \babelsublr{EPSG}).
  \item الفصل~2 (كوكب صالح للحياة): الزهرة من \babelsublr{Mariner~10} (\babelsublr{PIA23791})، \babelsublr{NASA/JPL-Caltech}، ملك عام؛ وجبل شارب كما صوّرته المركبة الجوالة \babelsublr{Curiosity} (\babelsublr{PIA16769})، \babelsublr{NASA/JPL-Caltech/MSSS}، ملك عام (عبر \babelsublr{Wikimedia Commons}).
  \item الفصل~2: بيانات الكواكب من \babelsublr{NASA fact sheets}؛ وحدود أطوار الماء من إصدارات \babelsublr{IAPWS} \babelsublr{SR1-86} و\babelsublr{R14-08}؛ وحدود النطاق الصالح للحياة من كوبارابو وآخرين (\babelsublr{Kopparapu and others}، 2013)؛ وبروكسيما~ب من \babelsublr{NASA Exoplanet Archive}.
  \item الفصل~3 (ميزانية الطاقة للأرض): برك ذوبان على الجليد البحري في المحيط المتجمد الشمالي (بعثة \babelsublr{ICESCAPE}، 2011)؛ صورة من \babelsublr{NASA Goddard Space Flight Center}، ملك عام (عبر \babelsublr{Wikimedia Commons}).
  \item الفصل~3: قيم ميزانية الطاقة من تقرير التقييم السادس للهيئة الحكومية الدولية المعنية بتغير المناخ (\babelsublr{IPCC Sixth Assessment Report, Working Group~I})، الشكل~7.2؛ وقيم الألبيدو من \babelsublr{University of Colorado Boulder} (\babelsublr{MOSAiC}) و\babelsublr{National Snow and Ice Data Center}؛ ونقل الحرارة من ترنبرث وكارون (\babelsublr{Trenberth and Caron}، 2001)؛ والإنتاج الأولي الصافي من فيلد وآخرين (\babelsublr{Field and others}، 1998)؛ والبيانات الكيميائية الحرارية من \babelsublr{NIST Chemistry WebBook}.
  \item الفصل~4 (لماذا تتكون السحب): سحب ركامية في طقس صحو، صورة التقطها رالف ف.~كريسغ (\babelsublr{Ralph F. Kresge})، \babelsublr{NOAA}؛ وسندان سحابة ركامية مزنية، \babelsublr{NOAA National Severe Storms Laboratory}؛ كلتاهما ملك عام (عبر \babelsublr{Wikimedia Commons}).
  \item الفصل~4: التشبع محسوب من إصدارات \babelsublr{IAPWS}؛ وأحجام القطرات من \babelsublr{NASA Jet Propulsion Laboratory}.
  \item الفصل~5 (العواصف والأعاصير المدارية): الإعصار إيزابيل كما صُوّر من محطة الفضاء الدولية، صورة التقطها مايك ترنشارد (\babelsublr{Mike Trenchard})، \babelsublr{NASA Johnson Space Center}؛ وإعصار قمعي في كانساس، صورة التقطها شون وو (\babelsublr{Sean Waugh})، \babelsublr{NOAA National Severe Storms Laboratory}؛ وبرق، صورة التقطها س.~كلارك (\babelsublr{C.~Clark})، \babelsublr{NOAA National Severe Storms Laboratory}؛ كلها ملك عام (عبر \babelsublr{Wikimedia Commons}).
  \item الفصل~5: مسارات الأعاصير من \babelsublr{IBTrACS} الإصدار~\babelsublr{4r01} (غاهتان وآخرون، \babelsublr{Gahtan and others}، 2024؛ ناب وآخرون، \babelsublr{Knapp and others}، 2010؛ \babelsublr{NOAA National Centers for Environmental Information})؛ وفئات الأعاصير من \babelsublr{NOAA National Hurricane Center}؛ وطاقة الأعاصير وشروط تكونها من \babelsublr{NOAA Atlantic Oceanographic and Meteorological Laboratory}؛ ومعدل البرق من كريستيان وآخرين (\babelsublr{Christian and others}، 2003)؛ ومعلومات الأعاصير القمعية من \babelsublr{NOAA Storm Prediction Center}.
  \item الفصل~6 (المحيط الطبقي): وردة أخذ العينات، صورة التقطها بابلو كليمنتي-كولون (\babelsublr{Pablo Clemente-Colon})، \babelsublr{NOAA Photo Library}؛ وعوّامة قياس تُطلق من سفينة أبحاث، \babelsublr{NOAA Global Ocean Monitoring and Observing Program}؛ والكلوروفيل في المحيطات 2002--2010، جيسي ألن وروبرت سيمون (\babelsublr{Jesse Allen and Robert Simmon})، \babelsublr{NASA Earth Observatory}؛ كلها ملك عام (عبر \babelsublr{Wikimedia Commons}).
  \item الفصل~6: مقاطع المحيط من \babelsublr{World Ocean Atlas 2023} (ريغان وآخرون، \babelsublr{Reagan and others}، 2024؛ \babelsublr{NOAA National Centers for Environmental Information})؛ وكثافة ماء البحر من معادلة الحالة الدولية (فوفونوف وميلارد، \babelsublr{Fofonoff and Millard}، 1983)؛ وحجم المحيطات من شاريت وسميث (\babelsublr{Charette and Smith}، 2010)؛ ومعلومات العوامات من \babelsublr{Argo Program Office}؛ والصعود المائي من \babelsublr{NOAA National Ocean Service} وكيسلر (\babelsublr{Kessler}، 2006).
  \item الفصل~7 (التربة، مورد بطيء التكوّن): مقطع لاتيريت على النيس في كيرالا، صورة التقطها فيرنر شيلمان (\babelsublr{Werner Schellmann})، \babelsublr{CC~BY-SA~2.5}؛ وصخر متحلل (سابروليت) في وسط مدغشقر، صورة التقطها \babelsublr{Diorit}، \babelsublr{CC~BY-SA~3.0}؛ ومقطع تربة، \babelsublr{USDA Natural Resources Conservation Service}، ملك عام (كلها عبر \babelsublr{Wikimedia Commons}).
  \item الفصل~7: معدلات إنتاج التربة وتعريتها من مونتغمري (\babelsublr{Montgomery}، 2007)؛ وصيغ المعادن من \babelsublr{Mineralogy Database} (\babelsublr{webmineral.com})؛ وسمك الصخر المتحلل من \babelsublr{USGS Ground Water Atlas of the United States}.
  \item الفصل~8 (أحواض التصريف والأنهار والفيضانات): هيدرولوجي يقيس تصريف مجرى ومحطة قياس، \babelsublr{U.S. Geological Survey}، ملك عام (عبر \babelsublr{Wikimedia Commons}).
  \item الفصل~8: الذروات السنوية والتصريف اليومي لنهر بوتوماك قرب واشنطن، \babelsublr{U.S. Geological Survey National Water Information System}، المحطة \babelsublr{01646500}؛ وتصريف المسيسيبي وحوضه من كاميرر (\babelsublr{Kammerer}، 1990، \babelsublr{USGS})؛ وتصريف الأمازون من كاليد وآخرين (\babelsublr{Callède and others}، 2010)؛ وإرشادات السلامة من مياه الفيضانات من \babelsublr{NOAA National Weather Service}.
  \item الفصل~9 (موارد المياه العذبة): بحر آرال عامي 2000 و2018، \babelsublr{NASA Earth Observatory} (\babelsublr{MODIS/Terra})؛ وحقول مروية قرب سبليت، كانساس، \babelsublr{NASA} (\babelsublr{ASTER/Terra})، ملك عام (عبر \babelsublr{Wikimedia Commons}).
  \item الفصل~9: مخزونات الماء في العالم عن شيكلومانوف (\babelsublr{Shiklomanov}، 1993)، \babelsublr{U.S. Geological Survey Water Science School}؛ ودورة الماء العالمية عن ترنبرث وآخرين (\babelsublr{Trenberth and others}، 2007)، كما نُشرت في ترنبرث وأسرار (\babelsublr{Trenberth and Asrar}، 2014)؛ وأعمار المياه الجوفية في الحوض النوبي من ستورتشيو وآخرين (\babelsublr{Sturchio and others}، 2004)؛ وتغيرات خزان السهول العليا الجوفي من ماكغواير (\babelsublr{McGuire}، 2017، \babelsublr{USGS})؛ وميزانية بحر آرال من غايبولاييف وآخرين (\babelsublr{Gaybullaev and others}، 2012) وميكلين (\babelsublr{Micklin}، 2007).
  \item الفصل~10 (من الانجراف القاري إلى توسع قاع المحيط): سفينة الحفر \babelsublr{\emph{JOIDES Resolution}}، \babelsublr{Ocean Drilling Program}، \babelsublr{Texas A\&M University}، ملك عام (عبر \babelsublr{Wikimedia Commons}).
  \item الفصل~10: سلم أزمنة القطبية لكاندي وكنت (\babelsublr{Cande and Kent}، 1995)؛ ونتائج الحفر في الرحلة~3 من \babelsublr{Deep Sea Drilling Project} (1970)؛ وعمر قاع المحيط من سيتون وآخرين (\babelsublr{Seton and others}، 2020)، \babelsublr{EarthByte}، الرخصة \babelsublr{CC~BY~4.0}.
  \item الفصل~11 (الموجات الزلزالية تكشف باطن الأرض): تركيب مرسمة زلازل مؤقتة في بورتوريكو؛ صورة التقطها توماس برات (\babelsublr{Thomas Pratt})، \babelsublr{U.S. Geological Survey}، ملك عام (عبر \babelsublr{Wikimedia Commons}).
  \item الفصل~11: النموذج المرجعي الأولي للأرض لدزيفونسكي وأندرسون (\babelsublr{Preliminary Reference Earth Model}، \babelsublr{Dziewonski and Anderson}، 1981)؛ وأزمنة الانتقال التجريبية للموجات \babelsublr{P} و\babelsublr{S} من \babelsublr{Australian Passive Seismic Server} (\babelsublr{AuSPASS}).
  \item الفصل~12 (الغلاف الصخري والغلاف اللدن): جبل جليدي، \babelsublr{NOAA}، ملك عام (\babelsublr{USGS Water Science School}).
  \item الفصل~12: التدفق الحراري من ديفيز وديفيز (\babelsublr{Davies and Davies}، 2010)، الرخصة \babelsublr{CC~BY~3.0}؛ وأقاليم التدفق الحراري من روي وبلاكويل وديكر (\babelsublr{Roy, Blackwell and Decker}، 1972)؛ وخط بدء انصهار الوشاح لهيرشمان (\babelsublr{Hirschmann}، 2000)؛ وسمك القشرة وكثافتها من كريستنسن وموني (\babelsublr{Christensen and Mooney}، 1995)؛ والارتفاعات من \babelsublr{ETOPO1} (\babelsublr{NOAA}).
  \item الفصل~13 (الصفائح وحركاتها): محطة لتحديد المواقع بالأقمار الصناعية في سلسلة جبال كاسكيد؛ صورة التقطها ديف شيرود (\babelsublr{Dave Sherrod})، \babelsublr{U.S. Geological Survey}، ملك عام (عبر \babelsublr{Wikimedia Commons}). وتضاريس مظللة لقاع شمال المحيط الهادئ، \babelsublr{National Geophysical Data Center}، \babelsublr{NOAA}، ملك عام (عبر \babelsublr{Wikimedia Commons}).
  \item الفصل~13: أقطاب دوران الصفائح \babelsublr{NNR-MORVEL56} لأرغوس وغوردون وديميتس (\babelsublr{Argus, Gordon and DeMets}، 2011)؛ ومواقع المحطات من \babelsublr{Nevada Geodetic Laboratory} (بليويت وهاموند وكريمر، \babelsublr{Blewitt, Hammond and Kreemer}، 2018)؛ وأعمار البراكين لكلاغ ودالريمبل (\babelsublr{Clague and Dalrymple}، 1987)، \babelsublr{U.S. Geological Survey}؛ وحدود الصفائح لبيرد (\babelsublr{Bird}، 2003)؛ واليابسة من \babelsublr{Natural Earth}.
  \item الفصل~14 (الظهرات المحيطية والغلاف الصخري المحيطي): لابة وسائدية على صدع غالاباغوس، \babelsublr{NOAA Okeanos Explorer Program}، ملك عام؛ ومدخنة سوداء على ظهرة خوان دي فوكا، \babelsublr{University of Washington} و\babelsublr{NOAA}، ملك عام؛ والغابرو، صورة التقطتها دارلا ساندرول (\babelsublr{Darla Sondrol}، \babelsublr{GeoDIL})، \babelsublr{CC0}؛ وبيريدوتيت في بازلت، صورة التقطها جيمس سانت~جون (\babelsublr{James St.~John})، \babelsublr{CC~BY~2.0} (كلها عبر \babelsublr{Wikimedia Commons}).
  \item الفصل~14: أعماق قاع المحيط من \babelsublr{ETOPO1} (\babelsublr{NOAA}) وأعماره من سيتون وآخرين (\babelsublr{Seton and others}، 2020)، \babelsublr{EarthByte}، الرخصة \babelsublr{CC~BY~4.0}؛ وقوانين العمق لبارسونز وسكلاتر (\babelsublr{Parsons and Sclater}، 1977)؛ وخط بدء انصهار الوشاح لهيرشمان (\babelsublr{Hirschmann}، 2000)؛ وأقطاب الدوران \babelsublr{NNR-MORVEL56} لأرغوس وغوردون وديميتس (\babelsublr{Argus, Gordon and DeMets}، 2011).
  \item الفصل~15 (نطاقات الاندساس): الشيست الأزرق والإكلوجيت، صورتان التقطهما جيمس سانت~جون (\babelsublr{James St.~John})، \babelsublr{CC~BY~2.0}؛ وأضرار تسونامي في كيسينوما، اليابان، 2011، صورة التقطها ديلان ماكورد (\babelsublr{Dylan McCord})، \babelsublr{U.S. Navy}، ملك عام (كلها عبر \babelsublr{Wikimedia Commons}).
  \item الفصل~15: الزلازل من فهرس الزلازل لهيئة المسح الجيولوجي الأمريكية (\babelsublr{USGS}، \babelsublr{ComCat})؛ وحدود الصفائح ومعدلات تقاربها لبيرد (\babelsublr{Bird}، 2003)؛ وأعمار قاع المحيط من سيتون وآخرين (\babelsublr{Seton and others}، 2020)، \babelsublr{EarthByte}، الرخصة \babelsublr{CC~BY~4.0}؛ وسرعات الصفائح \babelsublr{NNR-MORVEL56} لأرغوس وغوردون وديميتس (\babelsublr{Argus, Gordon and DeMets}، 2011)؛ واليابسة من \babelsublr{Natural Earth}.
  \item الفصل~16 (المعادن): الكوارتز والأورثوكلاز والبيوتيت والأوليفين، صور التقطها شانون هاينلي ودارلا ساندرول وديكستر بيركنز ونيسا يول (\babelsublr{Shannon Heinle, Darla Sondrol, Dexter Perkins, Nessa Eull})، \babelsublr{Geology Digital Image Library}، \babelsublr{CC0}؛ والبازلت في شريحة رقيقة، صورتان التقطهما \babelsublr{Chmee2}، \babelsublr{CC~BY-SA~3.0}؛ والكوارتز في غرانيت في شريحة رقيقة، صورة التقطها مايكل ك.~رايغل (\babelsublr{Michael C.~Rygel})، \babelsublr{CC~BY-SA~3.0} (كلها عبر \babelsublr{Wikimedia Commons}).
  \item الفصل~16: تركيب القشرة القارية لرودنيك وغاو (\babelsublr{Rudnick and Gao}، 2014)؛ والأوزان الذرية القياسية من \babelsublr{IUPAC/CIAAW}؛ وصيغ المعادن وكثافاتها من \babelsublr{webmineral.com}؛ وكثافات الصخور الغرانيتية في سييرا نيفادا من \babelsublr{U.S. Geological Survey} (\babelsublr{Open-File Report 91-570-A}).
  \item الفصل~17 (الصخور النارية): أعمدة البازلت في جسر العمالقة، صورة التقطها آدم بيشوب (\babelsublr{Adam Bishop})، \babelsublr{CC~BY-SA~3.0}؛ وإل كابيتان، يوسميتي، \babelsublr{U.S. Geological Survey}، \babelsublr{CC0}؛ والبازلت والأنديزيت والريوليت والديوريت والغرانيت والأوبسيديان وصخر بورفيري وشريحتان رقيقتان، صور التقطها نيسا يول وتويلا بيكر وجوزيف ه.~هارتمان ومايك ديفيس وديكستر بيركنز ودارلا ساندرول (\babelsublr{Nessa Eull, Twyla Baker, Joseph H.~Hartman, Mike Davis, Dexter Perkins, Darla Sondrol})، \babelsublr{Geology Digital Image Library}، \babelsublr{CC0} (كلها عبر \babelsublr{Wikimedia Commons}).
  \item الفصل~17: تبرد الجدد والسيلات عن ديلاني (\babelsublr{Delaney}، 1986)، \babelsublr{U.S. Geological Survey Open-File Report 86-444}؛ وتراكيب الصهارات وترتيب التبلور من ملاحظات س.~أ.~نيلسون (\babelsublr{S.~A.~Nelson}) في الجيولوجيا الطبيعية ومن س.~إيرل (\babelsublr{S.~Earle})، \babelsublr{\emph{Physical Geology}} (\babelsublr{CC~BY~4.0})؛ وسيرة بوين من تاريخ \babelsublr{Geophysical Laboratory} الذي نشرته \babelsublr{Carnegie Science}.
  \item الفصل~18 (جيولوجيا الفحم والنفط والغاز): نباتات أحفورية في سجيل غني بالكربون، والليغنيت، والفحم البيتوميني، والأنثراسيت، صور التقطها جيمس سانت~جون (\babelsublr{James St.~John})، \babelsublr{CC~BY~2.0}؛ وطبقة ليغنيت في رومانيا، \babelsublr{U.S. Geological Survey}، ملك عام (كلها عبر \babelsublr{Wikimedia Commons}).
  \item الفصل~18: نافذة النفط وتكوّن الفحم والمصائد من س.~إيرل (\babelsublr{S.~Earle})، \babelsublr{\emph{Physical Geology}} (\babelsublr{CC~BY~4.0})؛ ومحتوى الفحم من الكربون من \babelsublr{U.S. Energy Information Administration}؛ والنظام البترولي عن ماغون وشموكر (\babelsublr{Magoon and Schmoker}، \babelsublr{U.S. Geological Survey}، 2000)؛ وأحواض الصدوع والملح في غرب وسط إفريقيا من \babelsublr{U.S. Geological Survey Fact Sheet 2023-3028}.
  \item الفصل~19 (القشرة القارية): نيس أكاستا، صورة التقطها \babelsublr{Pedroalexandrade}، \babelsublr{CC~BY-SA~3.0}؛ وحبيبات زركون جاك هيلز، صور التقطها أ.~كافوسي وج.~فالي (\babelsublr{A.~Cavosie and J.~Valley})، ملك عام؛ ونيس وغرانيت الدرع الكندي، صورة التقطها كريس لايت (\babelsublr{Chris Light})، \babelsublr{CC~BY-SA~4.0} (كلها عبر \babelsublr{Wikimedia Commons}).
  \item الفصل~19: تركيب القشرة لرودنيك وغاو (\babelsublr{Rudnick and Gao}، 2014)؛ وسمك القشرة من كريستنسن وموني (\babelsublr{Christensen and Mooney}، 1995)، ووايت وماكنزي وأونيونز (\babelsublr{White, McKenzie and O'Nions}، 1992)، وكاربلس وآخرين (\babelsublr{Karplus and others}، 2019)؛ وامتداد القارات من كوغلي (\babelsublr{Cogley}، 1984)؛ وأعمار قاع المحيط من سيتون وآخرين (\babelsublr{Seton and others}، 2020)، \babelsublr{EarthByte}، الرخصة \babelsublr{CC~BY~4.0}.
  \item الفصل~20 (بناء الجبال): جبل إيفرست من المكوك الفضائي كولومبيا، \babelsublr{NASA}، ملك عام؛ وطبقات صخرية مطوية في خانق كالي غانداكي، صورة التقطها جان-ماري أولو (\babelsublr{Jean-Marie Hullot})، \babelsublr{CC~BY~2.0}؛ وزحف غلاروس، صورة التقطها كريستيان هاينه (\babelsublr{Christian Heine})، \babelsublr{CC~BY-SA~3.0}؛ ولابة وسائدية في أفيوليت سمائل، صورة التقطها ليونيل دوشوازيل (\babelsublr{Lionel Duchoiselle})، \babelsublr{CC~BY-SA~3.0} (كلها عبر \babelsublr{Wikimedia Commons}).
  \item الفصل~20: المواقع الماضية للهند ولهاسا من نماذج الصفائح لمولر وآخرين (\babelsublr{Müller and others}، 2019، 2022)، محسوبة بخدمة \babelsublr{GPlates Web Service} (\babelsublr{EarthByte})؛ والمرتفعات من \babelsublr{ETOPO1}؛ والزلازل من فهرس \babelsublr{USGS}؛ وحدود الصفائح لبيرد (\babelsublr{Bird}، 2003)؛ واليابسة من \babelsublr{Natural Earth}.
  \item الفصل~21 (الصخور المتحولة): الانفصام الأردوازي تحت المجهر، صورة التقطها \babelsublr{Geoloshe}، \babelsublr{CC0}؛ وشيست ميكائي غارنيتي من سيروس، صورة التقطها غرايم تشيرتشارد (\babelsublr{Graeme Churchard})، \babelsublr{CC~BY~2.0}؛ وشيست الكيانيت ونيس الكوارتز والبيوتيت، صورتان التقطهما جيمس سانت~جون (\babelsublr{James St.~John})، \babelsublr{CC~BY~2.0} (كلها عبر \babelsublr{Wikimedia Commons}).
  \item الفصل~21: حقول استقرار الكيانيت والأندالوسيت والسيليمانيت من هولداواي وموخوبادياي (\babelsublr{Holdaway and Mukhopadhyay}، 1993).
  \item الفصل~22 (الصهارية في نطاقات الاندساس): عمود ثوران جبل سانت هيلينز، 18~مايو 1980، صورة التقطها دونالد أ.~سوانسون (\babelsublr{Donald A. Swanson})، \babelsublr{USGS}، ملك عام؛ والغرانوديوريت والأنديزيت وسييرا نيفادا، صور التقطها جيمس سانت~جون (\babelsublr{James St.~John})، \babelsublr{CC~BY~2.0} (كلها عبر \babelsublr{Wikimedia Commons}).
  \item الفصل~22: الزلازل من فهرس \babelsublr{USGS}؛ والبراكين من \babelsublr{Smithsonian Global Volcanism Program}؛ وحدود الصفائح لبيرد (\babelsublr{Bird}، 2003)؛ واليابسة من \babelsublr{Natural Earth}؛ وخط بدء انصهار البيريدوتيت الجاف لهيرشمان (\babelsublr{Hirschmann}، 2000)، مع كثافات \babelsublr{PREM}.
  \item الفصل~23 (هدم الجبال): جبال الأبالاش من محطة الفضاء الدولية، صورة التقطها أحد رواد الفضاء (\babelsublr{ISS033-E-22378})، \babelsublr{NASA}، ملك عام؛ ورصيص من حصى الكوارتز، صورة التقطها جيمس سانت~جون (\babelsublr{James St.~John})، \babelsublr{CC~BY~2.0} (كلتاهما عبر \babelsublr{Wikimedia Commons}).
  \item الفصل~23: معدلات التجريد لبورتنغا وبيرمان (\babelsublr{Portenga and Bierman}، 2011)؛ والحمولات الرسوبية للمسيسيبي لميد ومودي (\babelsublr{Meade and Moody}، 2010)؛ وثابت التناقص لأنيرت (\babelsublr{Ahnert}، 1970)؛ وتلال فاردن، المنطقة المتخيَّلة في هذه السلسلة.
  \item الفصل~24 (الحرارة الداخلية للأرض): محطة نيسيافيتلير للطاقة الحرارية الأرضية، آيسلندا، صورة التقطها غريتار إيفارسون (\babelsublr{Gretar Ívarsson})، وحرّرها \babelsublr{Fir0002}، ملك عام (عبر \babelsublr{Wikimedia Commons}).
  \item الفصل~24: التدفق الحراري لديفيز وديفيز (\babelsublr{Davies and Davies}، 2010)؛ والحرارة الإشعاعية من قياس النيوترينوات الأرضية في تجربة كاملاند (2011) ومن الحرارة التي تطلقها اليوم النظائر الأربعة؛ وإنتاج الصخور للحرارة من فيلا وفرنانديث (\babelsublr{Vilà and Fernández}، 2010).
  \item الفصل~25 (التأريخ الإشعاعي وعمر الأرض): وليام طومسون، اللورد كلفن، 1906، صورة من \babelsublr{Russell and Sons}، لندن، \babelsublr{Smithsonian Institution Libraries}، بلا قيود معروفة على حقوق النشر؛ ونيزك كانيون ديابلو الحديدي، صورة التقطها جيفري نوتكين (\babelsublr{Geoffrey Notkin})، \babelsublr{CC~BY-SA~2.5} (كلتاهما عبر \babelsublr{Wikimedia Commons}).
  \item الفصل~25: نظائر الرصاص في النيازك لباترسون (\babelsublr{Patterson}، 1956)؛ وثابت تفكك الروبيديوم-87 الذي أوصى به \babelsublr{IUGS} و\babelsublr{IUPAC} (2022)؛ وتاريخ التقديرات من هولمز (\babelsublr{Holmes}، 1913، 1927)، وباترسون (\babelsublr{Patterson}، 1956)، و\babelsublr{U.S. Geological Survey}؛ وتلال فاردن، المنطقة المتخيَّلة في هذه السلسلة.
  \item الفصل~26 (الأرض الفتية وغلافها الجوي): تكوين حديد شريطي، خانق ديلز، صورة التقطها غرايم تشيرتشارد (\babelsublr{Graeme Churchard})، \babelsublr{CC~BY~2.0}؛ وستروماتوليتات في خليج القرش، صورة التقطها بول هاريسون (\babelsublr{Paul Harrison})، \babelsublr{CC~BY-SA~3.0}؛ وطبقات حمراء قرب سيدونا، صورة التقطتها ويندي من بنسلفانيا (\babelsublr{Wendy from Pennsylvania})، \babelsublr{CC~BY-SA~2.0} (كلها عبر \babelsublr{Wikimedia Commons}).
  \item الفصل~26: لمعان الشمس الفتية لغوف (\babelsublr{Gough}، 1981)؛ وحدود الدهور وألوانها من الجدول الزمني الطبقي الدولي (\babelsublr{International Chronostratigraphic Chart}).
  \item الفصل~27 (قراءة المناخات الماضية): فقاعات الهواء في جليد القارة القطبية الجنوبية، صورة من \babelsublr{CSIRO Atmospheric Research}، \babelsublr{CC~BY~3.0}؛ ومنخربات طافية، صورة التقطتها مارسي روبنسون (\babelsublr{Marci Robinson})، \babelsublr{USGS}، ملك عام؛ وحبوب لقاح، \babelsublr{Dartmouth Electron Microscope Facility}، ملك عام (كلها عبر \babelsublr{Wikimedia Commons}).
  \item الفصل~27: درجة الحرارة في دوم~C من \babelsublr{EPICA} (جوزيل وآخرون، \babelsublr{Jouzel and others}، 2007) وثاني أكسيد الكربون في العينات الجليدية من القارة القطبية الجنوبية (بيرايتر وآخرون، \babelsublr{Bereiter and others}، 2015) عبر \babelsublr{NOAA Paleoclimatology}؛ ومجموعة \babelsublr{LR04} (ليسيكي ورايمو، \babelsublr{Lisiecki and Raymo}، 2005)؛ وتضاريس \babelsublr{ETOPO1} (\babelsublr{NOAA})؛ و\babelsublr{NASA GISS} (\babelsublr{GISTEMP v4})؛ والفترات المرجعية من تقرير التقييم السادس للهيئة الحكومية الدولية المعنية بتغير المناخ (\babelsublr{IPCC Sixth Assessment Report, Working Group~I})، الفصل~2.
  \item الفصل~28 (تغير المناخ): مرصد ماونا لوا؛ صورة التقطها براين فاسل (\babelsublr{Brian Vasel})، \babelsublr{NOAA}، ملك عام (عبر \babelsublr{Wikimedia Commons}).
  \item الفصل~28: تشارلز ديفيد كيلينغ وجون تشين في ماونا لوا، 1997؛ صورة التقطها جون ميلر (\babelsublr{John Miller})، \babelsublr{NOAA}، ملك عام (عبر \babelsublr{Wikimedia Commons}).
  \item الفصل~28: البيانات من \babelsublr{NOAA Global Monitoring Laboratory} و\babelsublr{Scripps Institution of Oceanography} (ثاني أكسيد الكربون)، و\babelsublr{NASA GISS} (\babelsublr{GISTEMP v4})، و\babelsublr{NOAA Laboratory for Satellite Altimetry} (مستوى سطح البحر)، وبيرايتر وآخرين (\babelsublr{Bereiter and others}، 2015) عبر \babelsublr{NOAA Paleoclimatology} (العينات الجليدية)، وتقرير التقييم السادس للهيئة الحكومية الدولية المعنية بتغير المناخ (\babelsublr{IPCC Sixth Assessment Report, Working Group~I}) (التأثيرات الإشعاعية، والإسناد، والإسقاطات).
\end{itemize}
